\subsection{Objective and experimental design}

Task~7 compares a convolutional classifier (Subtask~7-1) with the requested
non-convolutional supervised regression system (Subtask~7-2).  The seven
classes are \texttt{0}, \texttt{4}, \texttt{7}, \texttt{8}, \texttt{A},
\texttt{D}, and \texttt{H}.  The implementation is in
\path{task_7/task_7.py}; all numerical results, split manifests, checkpoints,
confusion matrices, and tuning curves are saved in
\path{task_7/outputs/}.  The supplied Figure~1 image and the ten crops made in
Task~6 are used only after training for inference, so no target-row pixels can
leak into the learned parameters.

\subsection{Data split and leakage control}

There are 1,016 grayscale $128\times128$ PNG images in each class, giving
7,112 labelled samples.  The base random seed is 20261013; deterministic per-candidate and per-member
seeds are derived from it as specified below.  The
outer split is exactly 75\% training and 25\% test: 762 images per class are
assigned to the outer training pool and 254 images per class to the untouched
test set.  Ten percent of the outer training pool (76 images per class) is
held out as validation for model selection, leaving 686 images per class for
the tuning fit.  After selecting the input resolution and learning rate, the
final models are refitted on all 5,334 images in the 75\% outer training pool.

The file names contain repeated acquisition frames.  A random file-level
split would therefore allow identical images to occur on both sides of the
evaluation boundary.  The script hashes decoded grayscale pixels with
SHA-256, keeps every exact duplicate group in one partition, and solves a
small subset-sum problem per class to retain the requested 762/254 counts.
The resulting train/validation/test counts are 4,802/532/1,778; the three
pairwise exact-pixel digest intersections are all zero.  The complete audit is
in \path{task_7/outputs/split_summary.json} and the reproducible assignment of
every file is in \path{task_7/outputs/split_manifest.csv}.

\subsection{Common preprocessing}

Images are converted to one channel, resized with bilinear interpolation to a
candidate square side ($32$ or $48$ pixels), and divided by 255.  The mean and
standard deviation are fitted only on the tuning-train subset and the same
statistics are applied to validation and test samples:
\begin{equation}
  z = \frac{x-\mu_{\mathrm{train}}}{\sigma_{\mathrm{train}}}.
  \label{eq:task7-normalisation}
\end{equation}
No augmentation is used because the seven glyphs are already centred and
small geometric changes can alter the distinguishing strokes.  The selected
normalisation statistics are recorded in \path{task7_metrics.json}.

\subsection{Subtask 7-1: CNN classifier}

The CNN uses two $3\times3$ convolution--BatchNorm--ReLU layers with 32
channels, max pooling, two corresponding layers with 64 channels, and one
128-channel convolution followed by global average pooling.  A linear layer
maps the pooled feature to seven logits.  Cross-entropy is used for the
classification loss.  Global average pooling makes the classifier insensitive
to small translations while keeping the number of parameters at 140,071.
The forward path is
\[
  1\times s\times s\;\rightarrow\;32\;\rightarrow\;64\;\rightarrow\;128
  \rightarrow\mathrm{GAP}\rightarrow7,
\]
where $s$ is the selected input side.  BatchNorm layers are included after
each convolution~\cite{batchnorm}, and AdamW with weight decay $10^{-4}$
is used for fitting~\cite{adamw}.

\subsection{Subtask 7-2: MLP one-hot regression}

The non-CNN method is a fully connected multi-output regression model.  For a
class $c$, the target is the one-hot vector $y\in\{0,1\}^{7}$ and the network
outputs seven real-valued scores $f_\theta(x)$.  It is trained with mean
square error,
\begin{equation}
  \mathcal{L}_{\mathrm{reg}}(\theta)=
  \frac{1}{7}\left\|f_\theta(x)-y\right\|_2^2,
  \qquad \widehat c=\arg\max_j f_{\theta,j}(x).
  \label{eq:task7-regression}
\end{equation}
This explicitly implements the requested supervised regression formulation;
the MLP does not use convolution or a cross-entropy classification head.  Its
structure is
\[
  s^2\;\rightarrow\;512\;\rightarrow\;256\;\rightarrow\;7,
\]
with GELU activations, BatchNorm1d and Dropout(0.25) after both hidden linear
layers~\cite{batchnorm,dropout}.  The selected model has 659,463 parameters.  Three independently
seeded MLPs are fitted and their seven-dimensional regression outputs are
averaged before the argmax decision.  This ensemble reduces variance without
changing the non-CNN nature of the method.

\subsection{Hyperparameter and preprocessing investigation}

The validation grid contains both $32\times32$ and $48\times48$ resizing and
learning rates $10^{-3}$ and $3\times10^{-4}$.  All eight runs are listed in
\cref{tab:task7-tuning}, and their validation curves are shown in
\cref{fig:task7-tuning}.  Every tuning run uses a single network; the three-member
MLP ensemble is fitted only after its configuration is selected.  The complete
machine-readable grid is in \path{task_7/outputs/tuning_results.csv}.

\begin{table}[!htbp]
  \centering
  \caption{Complete single-network validation grid. Bold configurations are
  selected for final fitting. CNN loss is cross-entropy; MLP loss is one-hot
  MSE, so the two loss scales should not be compared directly.}
  \label{tab:task7-tuning}
  \setlength{\tabcolsep}{5pt}
  \begin{tabular}{llrrrr}
    \toprule
    Method & Input & Initial LR & Best epoch & Val. accuracy & Val. loss \\
    \midrule
    CNN & $32\times32$ & $10^{-3}$ & 18 & 0.9962 & 0.0157 \\
    CNN & $32\times32$ & $3\times10^{-4}$ & 7 & 0.9887 & 0.1413 \\
    \textbf{CNN} & $\mathbf{48\times48}$ & $\mathbf{10^{-3}}$ & \textbf{15} & \textbf{0.9981} & \textbf{0.0416} \\
    CNN & $48\times48$ & $3\times10^{-4}$ & 18 & 0.9868 & 0.0962 \\
    \midrule
    MLP regression & $32\times32$ & $10^{-3}$ & 18 & 1.0000 & 0.0055 \\
    \textbf{MLP regression} & $\mathbf{32\times32}$ & $\mathbf{3\times10^{-4}}$ & \textbf{17} & \textbf{1.0000} & \textbf{0.0079} \\
    MLP regression & $48\times48$ & $10^{-3}$ & 17 & 1.0000 & 0.0046 \\
    MLP regression & $48\times48$ & $3\times10^{-4}$ & 16 & 0.9981 & 0.0074 \\
    \bottomrule
  \end{tabular}
\end{table}

Configuration selection maximises validation accuracy.  An exact tie is
resolved first by the smaller input side and then by the smaller learning
rate.  Thus the MLP's three 100\% validation runs are tied; selecting
$32\times32$ and $3\times10^{-4}$ follows this deterministic rule, rather than
showing that this resolution is more accurate.  Within a run, the best epoch
maximises accuracy and uses the lower validation loss to break an accuracy tie.

For the CNN, increasing the side from 32 to 48 raises validation accuracy from
99.62\% to 99.81\% at initial learning rate $10^{-3}$, but lowers it from
98.87\% to 98.68\% at $3\times10^{-4}$.  The observed resolution effect
therefore depends on the learning rate.  Initial rate $10^{-3}$ outperforms
$3\times10^{-4}$ for both CNN resolutions within this budget; MLP differences
are small and three settings reach the validation ceiling.  These are
exploratory grid results, not isolated causal estimates: each configuration
is run once, and its seed also changes according to
\texttt{base + 100*side + int(100000*lr) + model\_offset}, with offsets 0 for
CNN and 1 for MLP.  Random initialisation and batch order can contribute to
the small differences; repeated runs with shared seed sets would be needed
to quantify sensitivity more reliably.

\begin{table}[!htbp]
  \centering
  \caption{Shared training settings and the final-fit schedule.}
  \label{tab:task7-configuration}
  \begin{tabularx}{\linewidth}{@{}lX@{}}
    \toprule
    Setting & Value \\
    \midrule
    Optimiser & AdamW; $\beta_1=0.9$, $\beta_2=0.999$, $\epsilon=10^{-8}$; weight decay $10^{-4}$ \\
    Batch sizes & 128 for fitting, 256 for validation/test; no dropped final batch \\
    Tuning budget & At most 18 epochs; early-stop patience 5 without a better accuracy/loss checkpoint \\
    Tuning LR schedule & ReduceLROnPlateau on validation accuracy; factor 0.5, patience 2, minimum initial LR/20 \\
    Final fitting & All 5,334 outer-training samples; 15 CNN epochs or 17 epochs per MLP member; cosine LR decay with $T_{\max}$ equal to that epoch count \\
    Final seeds & CNN 20268013; MLP members 20268014, 20268051, 20268088 \\
    Normalisation & At 32 px: $\mu=0.8240095$, $\sigma=0.3436231$; at 48 px: $\mu=0.8239762$, $\sigma=0.3558236$; fitted only on tuning-train \\
    Loader/reproducibility & 0 workers; deterministic cuDNN; CUDA selected if available \\
    \bottomrule
  \end{tabularx}
\end{table}

\begin{figure}[!htbp]
  \centering
  \includegraphics[width=\textwidth]{tuning_curves.png}
  \caption{Validation accuracy and loss for the resolution and learning-rate
  investigation.}
  \label{fig:task7-tuning}
\end{figure}

\subsection{Held-out results and efficiency}

The untouched 1,778-image test set gives the results in
\cref{tab:task7-results}.  Both methods generalise well, but the CNN is
0.62 percentage points more accurate and has a slightly higher macro-F1.  The
MLP ensemble has a larger parameter count because every pixel connects to the
first dense layer; nevertheless, its lower-resolution input makes batched
inference faster in this small experiment.  Training time includes the final
fit (one CNN or three MLP members) on the complete outer training pool.

\begin{table}[!htbp]
  \centering
  \caption{Final held-out test performance and measured efficiency.
  Parameter counts are per member; training time covers the complete ensemble.}
  \label{tab:task7-results}
  \setlength{\tabcolsep}{4pt}
  \begin{tabularx}{\linewidth}{@{}>{\raggedright\arraybackslash}Xrrrrrr@{}}
    \toprule
    Method & Accuracy & Macro-F1 & Params & Members & \shortstack{Train\\(s)} & \shortstack{Infer.\\(ms/img)} \\
    \midrule
    CNN ($48\times48$) & 0.9933 & 0.9932 & 140,071 & 1 & 4.73 & 0.0375 \\
    MLP reg. ensemble ($32\times32$) & 0.9871 & 0.9871 & 659,463 & 3 & 7.19 & 0.0142 \\
    \bottomrule
  \end{tabularx}
\end{table}

The MLP test regression MSE is 0.0084 before the argmax decision.  Its
confusion matrix, like the CNN matrix, is nearly diagonal; the largest residual
confusions involve the visually similar \texttt{A}/\texttt{D}/\texttt{H}
strokes.  The matrices are provided as report-ready figures in
\cref{fig:task7-confusion}.

\begin{figure}[!htbp]
  \centering
  \begin{subfigure}[t]{0.48\textwidth}
    \centering
    \includegraphics[width=\linewidth]{cnn_confusion_matrix.png}
    \caption{CNN}
  \end{subfigure}\hfill
  \begin{subfigure}[t]{0.48\textwidth}
    \centering
    \includegraphics[width=\linewidth]{mlp_confusion_matrix.png}
    \caption{MLP regression ensemble}
  \end{subfigure}
  \caption{Confusion matrices on the duplicate-isolated held-out test set.}
  \label{fig:task7-confusion}
\end{figure}

\subsection{Inference on the Figure-1 character row}

After the held-out evaluation, the saved checkpoints were applied to the ten
Task~6 character crops.  None of these crops enters the split manifest,
normalisation fitting, network optimisation, or hyperparameter selection.
The original grayscale baseline inverts the bright-on-dark photograph,
square-pads each crop, resizes it, and applies the training normalisation.
It gives only 10\% CNN character accuracy and 90\% MLP character accuracy
on the known ten-character reference row (\cref{tab:task7-target}).  This
large CNN gap from the clean-dataset test score motivates a separate
preprocessing diagnosis.

A single class-agnostic alignment rule was fixed before inspecting its new
predictions: use the Task~6 binary crop, keep its largest 8-connected
foreground component, crop the component's tight bounding box, invert to
black strokes on white, and resize this box to the median relative
foreground width and height measured on the 4,802 tuning-train images.
Those medians are 0.609375 and 0.7578125; source foreground is defined by
pixel value below 128.  The resized glyph is centred on a white square canvas
and then receives the original trained model's normalisation.  This rule
removes photograph texture and aligns glyph scale; it can also distort aspect
ratio, since width and height are matched separately.  No target labels,
model scores, or held-out test measurements select this rule or its parameters.
The grayscale baseline and fixed binary alignment are shown in
\cref{fig:task7-target-preprocessing}.

\begin{table}[!htbp]
  \centering
  \caption{Figure~1 predictions, with ten ordered positions made explicit.
  Accuracy is measured on this single reference row after inference.}
  \label{tab:task7-target}
  \setlength{\tabcolsep}{4pt}
  \begin{tabularx}{\linewidth}{@{}>{\raggedright\arraybackslash}X*{10}{c}r@{}}
    \toprule
    Position & 1 & 2 & 3 & 4 & 5 & 6 & 7 & 8 & 9 & 10 & Accuracy \\
    \midrule
    Reference & H & D & 4 & 4 & 7 & 8 & 0 & A & 0 & 0 & --- \\
    CNN, gray & A & A & A & A & A & A & A & A & A & A & 1/10 (10\%) \\
    MLP ensemble, gray & H & D & 4 & 4 & 7 & H & 0 & A & 0 & 0 & 9/10 (90\%) \\
    CNN, aligned binary & A & D & 4 & 4 & 7 & 8 & 0 & A & 8 & 0 & 8/10 (80\%) \\
    MLP ensemble, aligned binary & H & D & 4 & 4 & 7 & 8 & 0 & A & 0 & 0 & 10/10 (100\%) \\
    \bottomrule
  \end{tabularx}
\end{table}

\begin{figure}[!htbp]
  \centering
  \includegraphics[width=\linewidth]{target_preprocessing.png}
  \caption{The original inverted grayscale inputs (top) retain background
  texture. Fixed binary foreground and training-median geometry alignment
  (bottom) produce cleaner, centred inputs without fitting on Figure~1.}
  \label{fig:task7-target-preprocessing}
\end{figure}

The unchanged checkpoints predict \texttt{AD44780A80} with the aligned CNN
and \texttt{HD44780A00} with the aligned MLP ensemble.  The CNN still confuses
the first \texttt{H} with \texttt{A} and the ninth \texttt{0} with \texttt{8}.
These are exploratory results on one ten-character row, not a new held-out
benchmark or evidence of general target-domain accuracy.  The controlled
clean-dataset test results and measured training times above are unchanged.
The results support improving foreground preparation before adding ensemble
members; additional evidence would require separately acquired target-domain
samples.  Any subsequent domain-aware retraining should use extra labelled
training images or augmentations of the allowed training partition, while
keeping Figure~1 and the held-out test partition out of fitting and selection.
The saved \path{target_predictions.json} records both preprocessing paths,
per-character scores, geometry-fit provenance, and checkpoint files.

\subsection{Reproduction and environment}

The code uses Python and PyTorch; no scikit-learn or pretrained model is used.
The verified runtime contains Python 3.8, PyTorch 2.0.0+cu118,
torchvision 0.15.1+cu118, NumPy 1.24.2, Pillow 9.4.0, and Matplotlib 3.7.1,
running on an NVIDIA GeForce RTX 4080 SUPER.  The complete installation,
Tasks~1--6 dependencies, full and smoke experiment commands, and saved-model
\texttt{--inference-only} instructions are provided in
\cref{sec:reproduction-guide}.  Separate output directories preserve the
reported experiment during reproduction.  The supplied
\path{task_7/environment.yml} pins the required Python dependencies, while
\path{task_7/outputs/task7_metrics.json} records the selected settings,
held-out results, and both target preprocessing paths.

The implementation uses PyTorch modules for tensor operations and automatic
differentiation~\cite{pytorch}, Pillow for image decoding/resizing~\cite{pillow},
NumPy for split auditing and metrics~\cite{numpy}, and Matplotlib for
figures~\cite{matplotlib}.  These external libraries are
acknowledged here as required by the project brief; the data split, model
definitions, regression target, metric calculation, and experiment orchestration
are implemented in the submitted script.

\FloatBarrier
